
相容条件 7.3.13、7.3.14、7.3.17
的基本想法如下：
当 $t_i$ 趋于 $1$ 时，
多值截面成为沿遗忘映射的拉回；
当 $t_i$ 趋于 $0$ 时，
它与只在第 $i$ 个位置插入 $L$
所得的多值截面一致。

为表述相容条件 7.3.13，先引入一些记号。
把
\[
\vec a=(\vec a(1),\ldots,\vec a(|\vec a|))
\in\{1,\ldots,k\}^{|\vec a|}
\]
分成不交并
\[
\vec a=\vec a^1\sqcup\vec a^2.
\]
从 $\vec P^+$ 中删去
$\vec a^1(1),\ldots,
\vec a^1(|\vec a^1|)$
这些位置上的 $L$，得到 $\vec P'$。
定义 $\vec a^{2\prime}$ 为
\[
\vec a^{2\prime}(j)
=\vec a^2(j)
-\#\{i\mid\vec a^1(i)<\vec a^2(j)\}.
\]
注意，$\vec P'$ 中第
$\vec a^{2\prime}(j)$ 条链仍为 $L$。
再定义
\[
I^1_{\vec a^1,\vec a^2}\colon
[0,1]^{|\vec a^2|}\longrightarrow
[0,1]^{|\vec a|}
\]
如下：若
$I^1_{\vec a^1,\vec a^2}
(t_1,\ldots,t_{|\vec a^2|})
=(s_1,\ldots,s_{|\vec a|})$，则
\[
s_i=
\begin{cases}
1,&\vec a(i)\in\vec a^1,\\
t_j,&\vec a(i)=\vec a^2(j).
\end{cases}
\]
第二支的等号按位置编号理解
\translatornote{源书此处分段式第二支写
$\vec a(i)\in\vec a^2(j)$；
$\vec a^2(j)$ 是单个位置编号，
故依紧随的例 7.3.12 改用等号。
后文 $I^0$ 的同类分段式亦然。}。

\begin{example}{7.3.12}\label{ex:7-3-12}
若 $\vec a=\{1,3,5\}$、
$\vec a^1=\{1,5\}$、
$\vec a^2=\{3\}$、$k=6$，
则 $\vec P=(P_1,P_2,P_3)$、
$\vec a^{2\prime}=\{2\}$、
$\vec P'=(P_1,L,P_2,P_3)$、
$\vec P^+=(L,P_1,L,P_2,L,P_3)$，且
$I^1_{\vec a^1,\vec a^2}(t)=(1,t,1)$。
\end{example}

与 (7.3.7) 同样地定义遗忘映射
\[
\operatorname{forget}_{\vec a^1,\vec a^2}\colon
\mathcal M^{\mathrm{main}}_{k+1}(\beta;\vec P^+)
\longrightarrow
\mathcal M^{\mathrm{main}}_{k+1-|\vec a^1|}
(\beta;\vec P').
\]
它与仓西结构相容的方式
与引理~\ref{lem:7-3-8} 相同。
\begin{samepage}
令 $\mathfrak d'$ 为 $\vec P'$ 对应的装饰。
相应的多值截面是
\[
\mathfrak s^+_{\vec a^{2\prime},\mathfrak d',
k-|\vec a^1|,\beta,\vec P'}
\]
其定义域为
\[
[0,1]^{|\vec a^{2\prime}|}\times
\mathcal M^{\mathrm{main}}_{k+1-|\vec a^1|}
(\beta;\vec P')
\]
这里的第三个指标按定义域取值
\translatornote{源书此处及相容条件 7.3.13
把遗忘后截面的第三个指标仍写作 $k$；
但其定义域只有 $k-|\vec a^1|$ 个输入链，
故依模空间型别补全该指标。}。
\end{samepage}

\begin{statement}{相容条件}{7.3.13}\label{cond:7-3-13}
\liftIndex{I283}%
若
$\mathcal M^{\mathrm{main}}_{k+1-|\vec a^1|}
(\beta;\vec P')\ne\varnothing$，
则 $\mathfrak s^+_{\vec a,\mathfrak d^+,
k,\beta,\vec P^+}$ 在
$I^1_{\vec a^1,\vec a^2}\times\operatorname{id}$
的像上的限制，等于经
$\operatorname{forget}_{\vec a^1,\vec a^2}$
拉回的多值截面
$\mathfrak s^+_{\vec a^{2\prime},\mathfrak d',
k-|\vec a^1|,\beta,\vec P'}$。
\end{statement}

若模空间
$\mathcal M^{\mathrm{main}}_{k+1-|\vec a^1|}
(\beta;\vec P')$ 为空，
则按下述方式修改相容条件~\ref{cond:7-3-13}。

\begin{statement}{相容条件}{7.3.14}\label{cond:7-3-14}
先设 $k=2$、$\beta=\beta_0=0$，
并令 $\vec a=\vec a^1$ 为
$\{1\}$ 或 $\{2\}$。
也就是说，$\vec P=\{P\}$、
$\vec P'=(P)$，而
$\vec P^+=(L,P)$ 或 $(P,L)$
\translatornote{源书此处把遗忘后的
$\vec P'$ 写成空元组；但从 $\vec P^+$
只删去唯一的 $L$ 后，仍有输入链 $P$。}。
在这种情形下，
\[
\left.\mathfrak s^+_{\vec a,\mathfrak d^+,2,
\beta_0,\vec P^+}\right|_{
\{1\}\times\mathcal M^{\mathrm{main}}_3
(\beta_0;\vec P^+)}=0,
\]
而评价映射
$\operatorname{ev}_0\colon
[0,1]\times\mathcal M^{\mathrm{main}}_3
(\beta_0;\vec P^+)\to L$
的限制给出
\[
\begin{aligned}
\operatorname{ev}_{0*}^+\bigl(&
(\mathfrak s^+_{\vec a,\mathfrak d^+,2,
\beta_0,\vec P^+})^{-1}(0)\\
&\cap
(\{1\}\times\mathcal M^{\mathrm{main}}_3
(\beta_0;\vec P^+))\bigr)=P.
\end{aligned}
\]

在目标模空间
$\mathcal M^{\mathrm{main}}_{k+1-|\vec a^1|}
(\beta;\vec P')$ 为空的其他所有情形，
要求
\[
\begin{aligned}
\operatorname{ev}_{0*}^+\bigl(&
(\mathfrak s^+_{\vec a,\mathfrak d^+,k,
\beta,\vec P^+})^{-1}(0)\\
&\cap
(\{1\}\times\mathcal M^{\mathrm{main}}_{k+1}
(\beta;\vec P^+))\bigr)
\end{aligned}
\]
是退化链\translatornote{源书此处虽讨论任意
$\beta$ 的其他情形，却把截面和模空间的
下标都固定为 $\beta_0$；
依本分支的假设和命题 7.3.10 的量词改为 $\beta$。}。
\end{statement}

\begin{remark}{7.3.15}\label{rem:7-3-15}
在相容条件~\ref{cond:7-3-14} 的情形下，
说明
\[
\begin{aligned}
\operatorname{ev}_{0*}^+\bigl(&
(\mathfrak s^+_{\vec a,\mathfrak d^+,2,
\beta_0,\vec P^+})^{-1}(0)\\
&\cap
(\{0\}\times\mathcal M^{\mathrm{main}}_3
(\beta_0;\vec P^+))\bigr).
\end{aligned}
\]
若 $\mathfrak m_{2,\beta_0}$
是用两个微分同胚 $\varphi_0$、$\varphi_1$
定义的，则
$\mathfrak m_{2,\beta_0}(P,L)=\varphi_0(P)$、
$\mathfrak m_{2,\beta_0}(L,P)=\varphi_1(P)$。
因此，下文的相容条件 7.3.17 蕴含
\[
\begin{aligned}
\operatorname{ev}_{0*}^+\bigl(&
(\mathfrak s^+_{\vec a,\mathfrak d^+,2,
\beta_0,\vec P^+})^{-1}(0)\\
&\cap
(\{i\}\times\mathcal M^{\mathrm{main}}_3
(\beta_0;\vec P^+))\bigr)
=\varphi_i(P).
\end{aligned}
\]
这里的 $i=1$ 端需连同源书原样读法理解
\translatornote{此式按原书保留；
但取 $i=1$ 时，它给 $\varphi_1(P)$，
与相容条件 7.3.14 明写的 $t=1$ 端
评价链 $P$ 不一致；原书未假设
$\varphi_1(P)=P$，有待后文核对。}。
所以，相容条件~\ref{cond:7-3-14}
意味着
\[
P\longmapsto\operatorname{ev}_{0*}^+
\left(\mathcal M^{\mathrm{main}}_3
(\beta_0;\vec P^+)^{
\mathfrak s^+_{\vec a,\mathfrak d^+,2,
\beta_0,\vec P^+}}\right)
\]
是从恒等映射到
$(\varphi_i)_*$ 的一个“链同伦”。
\end{remark}

现在说明相容条件 7.3.17。
如上把 $\vec a$ 分为不交并
$\vec a^1\sqcup\vec a^2$，并定义
\[
I^0_{\vec a^1,\vec a^2}\colon
[0,1]^{|\vec a^2|}\longrightarrow
[0,1]^{|\vec a|}
\]
如下：若
$I^0_{\vec a^1,\vec a^2}
(t_1,\ldots,t_{|\vec a^2|})
=(s_1,\ldots,s_{|\vec a|})$，则
\[
s_i=
\begin{cases}
0,&\vec a(i)\in\vec a^1,\\
t_j,&\vec a(i)=\vec a^2(j).
\end{cases}
\]

\begin{example}{7.3.16}\label{ex:7-3-16}
若 $\vec a=\{1,3,5\}$、
$\vec a^1=\{1,5\}$、
$\vec a^2=\{3\}$、$k=6$，
则 $\vec P=(P_1,P_2,P_3)$、
$\vec P^+=(L,P_1,L,P_2,L,P_3)$，且
$I^0_{\vec a^1,\vec a^2}(t)=(0,t,0)$。
\end{example}

\begin{statement}{相容条件}{7.3.17}\label{cond:7-3-17}
\liftIndex{I284}%
$\mathfrak s^+_{\vec a,\mathfrak d^+,
k,\beta,\vec P^+}$ 在
$I^0_{\vec a^1,\vec a^2}\times\operatorname{id}$
的像上的限制，等于
$\mathfrak s^+_{\vec a^2,\mathfrak d^+,
k,\beta,\vec P^+}$。

当 $\vec a^2=\varnothing$
（即 $\vec a=\vec a^1$）时，
$\mathfrak s^+_{\vec a,\mathfrak d^+,
k,\beta,\vec P^+}$ 在
\[
\{0\}^{|\vec a|}\times
\mathcal M^{\mathrm{main},\mathfrak d^+}_{k+1}
(\beta;\vec P^+)
\]
上的限制为
$\mathfrak s_{\mathfrak d^+,k,\beta,\vec P^+}$。
\end{statement}
